\section{Tangent Vectors}

\subsection{Geometric tangent vectors}

\bdefn

    For $a\in\bR^n$ define the {\emphcolor geometric tangent space} at $a$ to be $\bR^n_a=\set a\times\bR^n$, endowed with the
    vector-space structure induced by the obvious bijection $\bR^n_a\cong\bR^n$.
    We denote a vector $(a,v)\in\bR^n_a$ by $v_a$ or $v|_a$.

\edefn

For a geometric tangent vector $v_a$, we can define the {\emphcolor directional derivative} at $a$ in the direction $v$:
$$ D_v|_a\colon C^\infty(\bR^n)\to\bR,\qquad D_v|_af = D_vf(a) = \eval{\frac d{dt}}_{t=0}f(a+tv) . $$
$D_v|_a$ is a linear map clearly, and it is also a {\emphcolor derivation}: by the product rule,
$$ D_v|_a(fg) = f(a)D_v|_a(g) + g(a)D_v|_a(f) . $$
Let $e_1,\dots,e_n$ form the standard basis for $\bR^n$, the chain rule gives us that for $v_a=v^ie_i|_a$ (Einstein summation),
then
$$ D_v|_af = v^i\frac{\partial f}{\partial x^i}(a) . $$
(Note that to use the summation convention on the right-hand side, we must view that the upper index in the denominator as a lower index.)
More succinctly,
$$ D_v|_af = v\cdot\nabla f(a) . $$

This leads to the following definition.

\bdefn

    A map $w\colon C^\infty(\bR^n)\to\bR$ is called a {\emphcolor derivation} at $a\in\bR^n$ if it is $\bR$-linear, and for all
    $f,g$:
    $$ w(fg) = f(a)\cdot w(g) + g(a)\cdot w(f) . $$
    Define $T_a\bR^n$ to be the set of all derivations at $a$.

\edefn

Clearly $T_a\bR^n$ has the natural structure of a linear space, via pointwise addition and multiplication, and all directional derivatives
are derivations.
More surprisingly, $T_a\bR^n$ is naturally isomorphic to $\bR^n_a$.

\blemm

    Let $w\in T_a\bR^n$ be a derivation at $a$.
    \benum
        \item $w$ is zero on constant functions.
        \item If $f(a)=g(a)=0$ then $w(fg)=0$.
    \eenum

\elemm

\bproof

    By linearity, it is sufficient to prove (1) on the constant function $f=1$.
    Notice that $f^2=f$, so
    $$ w(f) = w(f^2) = f(a)w(f) + f(a)w(f) = 2w(f), $$
    so $w(f)=0$, as desired.
    The second point is clear.
    \qed

\eproof

\bprop

    Let $a\in\bR^n$, then the map $\bR^n_a\to T_a\bR^n$ given by $v_a\mapsto D_v|_a$ is a linear isomorphism.

\eprop

\bproof

    Firstly, clearly this is linear (for a quick proof recall that $D_v|_a=v\cdot\nabla|_a$, and inner products are bilinear).
    To show that it is injective, suppose $D_v|_a=0$.
    In particular $D_v|_a(x^i)=0$ for all $i=1,\dots,n$ ($x^i$ denoting the $i$th projection).
    But then
    $$ D_v|_a(x^i) = v\cdot\nabla x^i(a) = v\cdot e_i = v^i , $$
    where $e_i$ is the $i$th basis vector.
    So $v^i=0$ for all $i$, meaning $v=0$.
    Thus the map has trivial kernel and is thus injective.

    To show surjectivity, notice that we would expect $w=D_v|_a$ where $v$ has components $w(x^i)$, as per the above computation.
    So let $v=w(x^i)e_i$, and we show that $D_v|_a=w$.
    Let $f\in C^\infty(\bR^n)$, and recall by Taylor's theorem we have
    $$ f(x) = f(a) + \sum_{i=1}^n\frac{\partial f}{\partial x^i}(a)(x^i-a^i) +
    \sum_{i,j=1}^n(x^i-a^i)(x^j-a^j)\int_0^1(1-t)\frac{\partial^2f}{\partial x^i\partial x^j}(a+t(x-a))\,dt . $$
    Notice that $x^i-a^i$ are all zero on $a$, so by the above lemma $w$ maps the right sum to zero ($f=x^i-a^i$ and $g$ is $x^j-a^j$
    times the integral).
    $w$ maps constants to zero, so
    $$ w(f) = \sum_{i=1}^n \frac{\partial f}{\partial x^i}(a)\cdot w(x^i) = v\cdot\nabla f(a) = D_v|_a(f) , $$
    as desired.
    \qed

\eproof

\bcoro

    For any $a\in\bR^n$, the $n$ derivations
    $$ \eval{\frac\partial{\partial x^1}}_a,\dots,\eval{\frac\partial{\partial x^n}}_a,\qquad \eval{\frac\partial{\partial x^i}}_a(f) = \frac{\partial f}{\partial x^i}(a) $$
    form a basis for $T_a\bR^n$.

\ecoro

This is simply since these are the derivations corresponding to the standard basis under the natural isomorphism from the above result.

\subsection{Tangent vectors on manifolds}

In Euclidean space we defined tangent vectors via a geometric construction, and showed that it was natural isomorphic to a more algebraic definition.
For general smooth manifolds we do not have the benefit of a geometric construction, and so we turn to the algebraic definition as above.

\bdefn

    Let $M$ be a smooth manifold (with or without boundary) and $p\in M$.
    A {\emphcolor derivation} at $p$ is a linear map $v\colon C^\infty(M)\to\bR$ satisfying
    $$ v(fg) = f(p)v(g) + g(p)v(f) $$
    for all $f,g$.
    Denote the space of derivations at $p$ by $T_pM$: this is a linear space as before.

\edefn

As before, derivations map constant functions to zero and if $p$ is a zero of both $f$ and $g$ then $v(fg)=0$.

\bdefn

    The tangent space construction is functorial: if $M,N$ are smooth manifolds, $p\in M$, and $F\colon M\to N$ is a smooth map,
    we define $F$'s {\emphcolor differential} at $p$ to be the map
    $$ dF_p\colon T_pM\to T_{F(p)}N $$
    defined by mapping a derivation $v\in T_pM$ to the derivation $dF_p(v)\in T_{Fp}N$ given by $dF_p(v)(f)=v(f\circ F)$.

\edefn

$dF_p$ is well-defined, since $f\circ F$ is a composition of smooth maps $M\to N\to\bR$ and thus is smooth.
$dF_p(v)$ is a derivation at $F(p)$ because clearly it is linear, and for $f,g\in C^\infty(N)$:
$$ \multline{
    dF_p(v)(fg) = v((fg)\circ F) = v((f\circ F)(g\circ F)) =\cr
    (f\circ F)(p)v(g\circ F) + (g\circ F)(p)v(f\circ F) = f(F(p))dF_p(v)(f) + g(F(p))dF_p(v)(g)
} $$
as desired.
$dF_p$ is also clearly linear.

To show functorality we need to show that it preserves identities and composition.
Clearly it does map $\id_M$ to $\id_{T_pM}$, and if $M\xto FN\xto GP$ are smooth then
$$ d(G\circ F)_p(v)\colon f\mapsto v(f\circ G\circ F), d(G)_{Fp}\circ d(F)_p(v)\colon f\mapsto d(F_p)(v)(f\circ G) = v(f\circ G\circ F) , $$
as desired.
In particular, for a diffeomorphism $F$, $dF_p\colon T_pM\to T_{Fp}M$ is an isomorphism with inverse $d(F^{-1})_{Fp}$.

\bprop

    Let $M$ be smooth and $p\in M$.
    If $f,g\in C^\infty(M)$ agree on some neighborhood of $p$, then $v(f)=v(g)$ for all derivations $v\in T_pM$.

\eprop

\bproof

    Setting $h=f-g$ and using linearity, it is sufficient to show that $v$ is zero on maps which are zero in a neighborhood of $p$.
    Since $h$ is zero in a neighborhood $U$ of $p$, $\supp h\subseteq M-U\subseteq M-p$, and thus we can take $\psi\in C^\infty(M)$ to be
    a bump function which is $1$ on $\supp h$ and supported in $M-p$.
    So $\psi h=h$ always, since $\psi=1$ when $h$ is nonzero.
    Thus $v(\psi h)=v(h)$, and since $\psi(p)=h(p)=0$, $v(\psi h)=0$ as desired.
    \qed

\eproof

Thus $T_pM$ is determined by the behavior of maps in neighborhoods of $p$ (this will lead to an equivalent definition of $T_pM$ as we will see soon).
In particular, we have the following:

\bprop

    Let $U\subseteq M$ be an open submanifold and let $\iota\colon U\to M$ be the inclusion map.
    Then for $p\in U$, the differential $d\iota_p\colon T_pU\to T_pM$ is an isomorphism.

\eprop

\bproof

    To show injectivity, suppose $d\iota_p(v)=0$ for $v\in T_pU$.
    Now let $f\in C^\infty(U)$ and let $p\in B\subseteq U$ be a neighborhood such that $\cl B\subseteq U$.
    Then we can extend $f$ to $\tilde f\colon M\to U$ such that $\tilde f$ agrees with $f$ on $\cl B$.
    By the above proposition, since $\tilde f$ and $f$ agree on the neighborhood $B$ of $p$, $v(f)=v(\tilde f|_U)$.
    Restriction is just composition with $\iota$, i.e.\ we have
    $$ v(f) = v(\tilde f|_U) = v(\tilde f\circ\iota) = d\iota_p(v)(\tilde f) = 0 $$
    since $d\iota_p(v)=0$.
    Thus $v(f)=0$ for all $f\in C^\infty(U)$, so $v=0$ as desired.
    Thus $d\iota_p$ is injective.

    To show surjectivity, for $w\in T_pM$ define $v\in T_pU$ by $v(f)=w(\tilde f)$ for the extension above.
    This is well-defined since any extension agree on a neighborhood of $p$.
    This is a valuation, as $\tilde f+\tilde g$, $c\tilde f$, and $\tilde f\cdot\tilde g$ are extensions of $f+g$, $cf$, $fg$ respectively.
    Thus linearity follows and
    $$ v(fg) = w(\tilde f\tilde g) = f(p)w(\tilde f) + g(p)w(\tilde g) = f(p)v(f) + g(p)v(g) $$
    as desired.
    Finally since $f\in C^\infty(M)$ is an extension of $f\circ\iota$, $d\iota_p(v)(f)=v(f\circ\iota)=w(f)$, showing surjectivity.
    \qed

\eproof

\bprop

    Let $M$ be a smooth manifold without boundary of dimension $n$, and $p\in M$.
    Then $T_pM$ is an $n$-dimensional real vector space.

\eprop

\bproof

    Let $(U,\phi)$ be a smooth chart containing $p$ to an open ball $B\subseteq\bR^n$.
    In particular $\phi$ forms a diffeomorphism $U\cong B$, which is in turn diffeomorphic to $\bR^n$, and the above proposition shows us that
    $$ T_pM \cong T_pU \cong T_{\phi p}\hat U \cong T_{\phi p}\bR^n , $$
    which is $n$-dimensional, as we showed.
    \qed

\eproof

What about manifolds with boundary?
For an interior point $p\in\Int M$ we have $T_pM\cong T_p\Int M$ which is $n$-dimensional (as $\Int M$ has no boundary).
So we need to consider only boundary points.
In particular we need to consider $T_a\bH^n$ for $a\in\partial\bH^n$.

\blemm

    Let $\iota\colon\bH^n\to\bR^n$ be the inclusion map.
    Then for any $a\in\partial\bH^n$, $d\iota_a\colon T_a\bH^n\to T_a\bR^n$ is an isomorphism.

\elemm

\bproof

    To show injectivity suppose $d\iota_a(v)=0$ for $v\in T_a\bH^n$.
    Then for $f\in C^\infty(\bH^n)$ we can extend $f$ to $\tilde f\in C^\infty(\bR^n)$ (since $\bH^n\subseteq\bR^n$ is closed).
    This means that $\tilde f\circ\iota=f$ and so
    $$ 0 = d\iota_a(v)(\tilde f) = v(\tilde f\circ\iota) = v(f) , $$
    meaning $v=0$, so $d\iota_a$ is injective.

    Suppose $w\in T_a\bR^n$, then define $v\in T_a\bH^n$ by $v(f)=w(\tilde f)$ for any extension $\tilde f$ of $f$.
    Now it is not immediately clear that this is well-defined, since two extensions may not agree on a neighborhood of $p\in\bR^n$.
    But we know that $w=w^i\partial/\partial x^i$ with respect to the standard basis of $T_a\bR^n$, and so
    $$ v(f) = w^i\frac{\partial\tilde f}{\partial x^i}(a) . $$
    By continuity, the partial derivatives of $\tilde f$ are determined by those of $f$ in a neighborhood of $a\in\bH^n$.
    So for any extension, the partial derivatives at $a$ are the same, so this is indeed well-defined.
    The proof that $v$ is a derivation and $d\iota_a(v)=w$ is the same as before.
    \qed

\eproof

\bprop

    Let $M$ be a manifold with or without boundary of dimension $n$.
    Then for each $p\in M$, $T_pM$ is $n$-dimensional.

\eprop

\bproof

    This is true as explained for interior points.
    For a boundary point $p\in\partial M$, let $(U,\phi)$ be a coordinate half-ball, $\phi$ being a diffeomorphism to $B\subseteq\bH^n$.
    Then $B$ is in turn diffeomorphic to $\bH^n$ and we have
    $$ T_pM \cong T_pU \cong T_pB \cong T_p\bH^n \cong T_p\bR^n , $$
    which is $n$-dimensional.
    \qed

\eproof

For a Euclidean space $V\cong\bR^n$, we can naturally identify $T_pV$ with $V$ itself, as we did with $\bR^n$.
The construction is the same: map $v\in V$ to the derivation
$$ D_v|_p(f) = \eval{\frac d{dt}}_{t=0}f(a+tv) . $$
Then the map $V\to T_pV$ given by $v\mapsto D_v|_p$ is a natural isomorphism.
The argument for it being an isomorphism is the same as before.
Naturality just means that it commutes with linear maps in the sense that if $L\colon V\to W$ is linear, the diagram below commutes

\centerline{\drawdiagram{
    $V$&$T_pV$\cr
    $W$&$T_{Lp}W$
}{
    \diagarrow{from={1,1}, to={1,2}, text=$\cong$, y distance=-.25cm}
    \diagarrow{from={2,1}, to={2,2}, text=$\cong$, y distance=.25cm}
    \diagarrow{from={1,1}, to={2,1}, text=$L$, x distance=-.25cm}
    \diagarrow{from={1,2}, to={2,2}, text=$d_pL$, x distance=.25cm}
}}

In particular, for an open submanifold of a vector space, its tangent space is naturally isomorphic to that vector space.
For example, $\gl_n(\bR)\subseteq M_n(\bR)$ is an open submanifold and as such $T_X\gl_n(\bR)\cong M_n(\bR)$ for every $X\in\gl_n(\bR)$.

\bprop

    Let $M_1,\dots,M_n$ be smooth manifolds (of which at most one has boundary) and let $M=M_1\times\cdots\times M_n$ be their direct product.
    Let $\pi_i\colon M\to M_i$ be the projection maps.
    Then for every $p=(p_1,\dots,p_n)\in M$ the differentials assemble into an isomorphism,
    $$ (d\pi_1,\dots,d\pi_n)\colon T_pM \to \bigoplus_{i=1}^n T_{p_i}M_i . $$

\eprop

\bproof

    Let $\alpha=(d\pi_1,\dots,d\pi_n)$, and define the inclusion maps $\iota_i\colon M_i\to M$ which maps $\iota_i(q)$ to be the point in $M$
    whose $i$th coordinate is $q$ and for $i\neq j$ the $j$th coordinate is $p_j$.
    $\iota_i$ is smooth which gives us maps $(d\iota_i)_{p_i}\colon T_{p_i}M_i\to T_pM$.
    Furthermore $\pi_i\circ\iota_i=\id_{M_i}$ and $\pi_j\circ\iota_i$ is constant, so $d(\pi_j\circ\iota_i)=\delta_{ij}$.

    Further define $\beta\colon\bigoplus_{i=1}^nT_{p_i}M_i\to T_pM$ by $\beta(v_1,\dots,v_n)=\sum_{i=1}^n(d\iota_i)_{p_i}(v_i)$.
    Then
    $$ \alpha\circ\beta(v_1,\dots,v_n) = \sum_{i=1}^n\alpha(d\iota_i)_{p_i}(v_i) , $$
    and by functorality
    $$ = \sum_{i=1}^n(d(\pi_1\iota_i)(v_i),\dots,d(\pi_n\iota_i)(v_i)) = \sum_{i=1}^n (0,\dots,v_i,\dots,0) = (v_1,\dots,v_n) . $$
    So $\alpha\circ\beta$ is the identity, and the dimensions of $T_pM$ and $\bigoplus T_{p_i}M_i$ are equal, so they are isomorphisms.
    \qed

\eproof

\subsection{Alternative definitions}

There exist many equivalent definitions of the tangent space of a manifold at a point.

\bdefn

    Let $M$ be a smooth manifold and $p\in M$.
    Consider then the set of all smooth functions $f\colon U\to\bR$ where $p\in U\subseteq M$ is an open neighborhood of $p$; denote this data
    by $(f,U)$.
    We define an equivalence of these functions, where $(f,U)$ and $(g,V)$ are identified if they are equal on a neighborhood of $p$.
    Let $C_p^\infty(M)$ denote the set of all these equivalence classes, called {\emphcolor germs}, we denote the germ of a function $f$ at $p$ by $[f]_p$.

    Note that the space of germs can naturally be endowed with a real algebra structure.
    A {\emphcolor derivation} of $C^\infty_p(M)$ is a linear map $v\colon C^\infty_p(M)\to\bR$ satisfying:
    $$ v[fg]_p = f(p)v[g]_p + g(p)v[f]_p . $$
    We can then define the tangent space of $M$ at $p$ to be the space of all such derivations, which can be given a natural real vector space
    structure.

\edefn

If we denote by $D_pM$ the space of all derivations of $C^\infty_p(M)$, it is not hard to see that $D_pM$ and $T_pM$ are naturally isomorphic.
A derivation of $C^\infty(M)$ defines a derivation of $C^\infty_p(M)$ and vice versa in a unique way.
The benefit of this definition via germs is that it is local, and highlights the locality of the tangent space.
In analytic manifolds (i.e.\ $C^k$ ones), this is definition gives the desired properties of a tangent space, as the global definition does not
reduce to a local one, since bump functions may not exist.

A more geometric definition can also be given.

\bdefn

    Let $M$ be a smooth manifold and $p\in M$.
    Consider the collection of all smooth curves centered at $p$: all smooth maps $\gamma\colon J\to M$ where $J\subseteq\bR$ is an open interval
    containing $0$ and $\gamma(0)=p$.
    Again we define an equivalence on such curves: consider $\gamma_1,\gamma_2$ equivalent if for every open neighborhood $p\in U\subseteq M$ and every
    smooth $f\colon U\to\bR$, $(f\circ\gamma_1)'(0)=(f\circ\gamma_2)'(0)$.
    Let $\V_p(M)$ be the collection of all these curves.

\edefn

$\V_p$ does not have an obvious linear structure, but we can give it one, and it is naturally isomorphic to $T_pM$.
We can also naturally define the differential of a smooth map $F\colon M\to N$ to be
$$ dF_p\colon\V_p(M)\to\V_{F(p)}(N),\qquad [\gamma]\mapsto[F\circ\gamma] . $$
This is well-defined: if $\gamma_1\sim\gamma_2$ then for $f\colon V\to\bR$ smooth in a neighborhood of $F(p)$, $f\circ F$ is smooth in a neighborhood
of $p$, and thus $(f\circ F\circ\gamma_1)'(0)=(f\circ F\circ\gamma_2)'(0)$ as desired.

\bprop

    Let $M$ be a smooth manifold, with or without boundary, and $p\in M$.
    Then $\Psi\colon\V_p(M)\to T_p(M)$ defined by $\Psi[\gamma](f)=(f\circ\gamma)'(0)$ is a bijection.

\eprop

\bproof

    Clearly $\Psi$ is well-defined, and it maps to derivations: clearly $\Psi[\gamma]$ is linear, and it is a derivation:
    $$ \Psi[\gamma](fg) = ((f\circ\gamma)\cdot(g\circ\gamma))'(0) = f(\gamma(0))\cdot(g\circ\gamma)'(0) + g(\gamma(0))\cdot(f\circ\gamma)'(0) =
    f(p)\Psi[\gamma](g) + g(p)\Psi[\gamma](f) . $$
    $\Psi$ is also injective; we can view this via the equivalence of derivations on $C^\infty_p(M)$ and $C^\infty(M)$.
    If $\Psi[\gamma_1]=\Psi[\gamma_2]$, then they define the same derivation on $C^\infty_p(M)$ and this clearly shows that $\gamma_1\sim\gamma_2$.

    Now let $(U,\phi)$ be a coordinate ball centered at $p$, so that $T_pM\cong T_pU\cong T_0B\cong\bR^n$.
    For $v\in\bR^n$, consider its tangent vector in $T_0B$ given by $D_v|_0$, which we can pull back via $d\phi^{-1}_0$ to give the tangent vector
    $$ d\phi^{-1}_0(D_v|_0) \in T_pU,\qquad d\phi^{-1}_0(D_v|_0)(f) = D_v|_0(f\circ\phi^{-1}) . $$
    This corresponds naturally to a tangent vector in $T_pM$ (which maps $f\in C^\infty(M)$ to its restriction to $U$ and then via the above formula).
    Define the curve
    $$ \gamma_v(t) = \phi^{-1}(tv) $$
    on an open interval $(-\epsilon,\epsilon)$ so that $tv\in U$ for all $t\in(-\epsilon,\epsilon)$.
    Then
    $$ \Psi[\gamma_v](f) = (f\circ\gamma_v)'(0) = \eval{\frac d{dt}}_0(f\circ\phi^{-1}(t\cdot v)) = D_v|_0(f\circ\phi^{-1}) , $$
    meaning $\Psi[\gamma_v]=d\phi^{-1}_0(D_v|_0)$, so that $\Psi$ is surjective as desired.
    \qed

\eproof

$\Psi$ can then be pulled back to give $\V_p(M)$ a vector space structure naturally isomorphic to $T_pM$.

This definition of the tangent space allows us to simply formalize the concept of a {\emphcolor velocity vector}.
Let $\gamma\colon J\to M$ be a smooth curve, then its velocity at $0$ is its equivalence class in $\V_{\gamma(0)}(M)$.
Its velocity at any $t_0\in J$ is the equivalence class of $\gamma_{t_0}(t)=\gamma(t_0+t)$ in $\V_{\gamma(t_0)}(M)$.

\subsection{Computations in coordinates}

Let $M$ be a smooth manifold without boundary and $(U,\phi)$ be a coordinate chart, so that
$$ d\phi_p\colon T_pM \to T_{\phi(p)}\bR^n $$
is an isomorphism for $p\in U$.
We know that the partial derivatives evaluated at $\phi(p)$, $\eval{\partial/\partial x^i}_{\phi(p)}$ for $1\leq i\leq n$, forms a basis for $T_{\phi(p)}\bR^n$.
Then their preimages under $d\phi_p$ forms a basis for $T_pM$.
Let us denote by $\ievpdv{x^i}_p$ these preimages, i.e.
$$ \evpdv{x^i}_p = (d\phi_p)^{-1}\parens{\evpdv{x^i}_{\phi(p)}} = d\phi^{-1}_{\phi(p)}\parens{\evpdv{x^i}_{\phi(p)}} . $$
Unwinding definitions, this valuation acts by
$$ \evpdv{x^i}_pf = \evpdv{x^i}_{\phi(p)}(f\circ\phi^{-1}) = \pdvof{\hat f}{\partial x^i}(\hat p) , $$
where $\hat f=f\circ\phi^{-1}$ is the coordinate representation of $f$, and $\hat p=(p^1,\dots,p^n)=\phi(p)$ is the coordinate representation of $p$.

When $M$ is a manifold with boundary, we need only consider the case of boundary points: for interior points we take interior charts.
For a boundary point $p\in\partial M$, the differential becomes an isomorphism
$$ d\phi_p\colon T_pM \to T_{\phi(p)}\bH^n , $$
and we can compose this with the differential $d\iota_{\phi(p)}\colon T_{\phi(p)}\bH^n\to T_{\phi(p)}\bR^n$.
The preimage of $\ipdv{x^i}_a$ under $d\iota_a$ is the one-sided partial derivative.
Then the preimages of $\ievpdv{x^i}_{\phi(p)}$ of $d(\iota\phi)_p$ gives us $\ievpdv{x^i}_p$ defined by
$$ \evpdv{x^i}_pf = \pdvof{\hat f}{\partial x^i}(\hat p) , $$
where the partial derivative is one-sided.

So $\evpdv{x^i}_p$ for $1\leq i\leq n$ forms a basis for $T_pM$, so any valuation $v\in T_pM$ can be uniquely written as
$$ v = v^i\evpdv{x^i}_p . $$
The basis $\evpdv{x^i}_p$ is called a {\emphcolor coordinate basis} for $T_pM$, and $(v^1,\dots,v^n)$ the {\emphcolor components} of $v$.
Notice that if we let $\pi_i\colon\bR^n\to\bR$ be the $i$th projection,
$$ v(\pi_i\circ\phi) = \parens{v^j\evpdv{x^j}_p}(\pi_i\circ\phi) = v^j\pdvof{x^i}{x^j}(\hat p) = v^i . $$
Oftentimes $\pi_i\circ\phi$ is identified with $x^i$, so we can write this concisely as
$$ v^i = v(x^i) , $$
giving an easy method of computing the components of $v$ relative to any coordinate basis.

\bprop

    Let $p\in M$ be a point in a smooth manifold.
    Then $\ievpdv{x^1}_p,\dots,\ievpdv{x^n}_p$ form a basis for $T_pM$, and for $v\in T_pM$ we have
    $$ v = v^i\evpdv{x^i}_p,\qquad\hbox{where }v^i=v(x^i) . $$

\eprop

Now how can we represent the differential of a smooth map $F\colon M\to N$ in coordinates?
To do so, we must compute the components of $(dF_p)(\ievpdv{x^i}_p)$ in $N$.

First suppose that $F\colon U\to V$ is a map from open Euclidean subsets $U\subseteq\bR^n,V\subseteq\bR^m$, and let $p\in U$.
Let $(x^1,\dots,x^n)$ denote the coordinates in the domain and $(y^1,\dots,y^m)$ in the codomain.
Then let us see how $dF_p$ maps the coordinate basis:
$$ dF_p\parens{\evpdv{x^i}_p}f = \evpdv{x^i}_p(f\circ F) = \pdvof{f}{y^j}(F(p))\pdvof{F^j}{x^i}(p) = \parens{\pdvof{F^j}{x^i}(p)\evpdv{y^j}_{F(p)}}(f) , $$
where the second equality is due to the chain rule.
Thus
$$ dF_p\parens{\evpdv{x^i}_p} = \pdvof{F^j}{x^i}(p)\cdot\evpdv{y^j}_{F(p)} , $$
meaning that the representation of $F$ relative to the coordinate bases is the matrix
$$ \pmatrix{\pdvof{F^1}{x^1}(p) & \cdots & \pdvof{F^1}{x^n}(p)\cr \vdots & \ddots & \vdots\cr \pdvof{F^n}{x^1}(p) & \cdots & \pdvof{F^n}{x^n}(p)} , $$
the {\emphcolor Jacobian} of $F$ at $p$.

Now in general suppose $F\colon M\to N$ is a smooth map between smooth manifolds, with $p\in M$.
Let $(U,\phi)$ be a chart in $M$ containing $p$, and $(V,\psi)$ a chart in $N$ containing $F(p)$.
Then let $\hat F=\psi\circ F\circ\phi^{-1}\colon\phi(U\cap F^{-1}(V))\to\psi(V)$ be its coordinate representation, $\hat p=\phi(p)$ be the coordinate representation
of $p$.
Then since $F\circ\phi^{-1}=\psi^{-1}\circ\hat F$,
$$ dF_p\parens{\evpdv{x^i}_p} = dF_p\parens{d(\phi^{-1})_{\hat p}\evpdv{x^i}_{\hat p}} = d(\psi^{-1})_{\hat F(\hat p)}\parens{d\hat F_p\parens{\evpdv{x^i}_{\hat p}}} . $$
We just computed the representation of $\hat F$:
$$ = d(\psi^{-1}_{\hat F(\hat p)})\parens{\pdvof{\hat F^j}{x^i}(\hat p)\evpdv{y^j}_{\hat F(\hat p)}} = \pdvof{\hat F^j}{x^i}\evpdv{y^j}_{F(p)} , $$
where the last equality is due to the definition of $\ievpdv{y^j}_{F(p)}$.
Thus the representation of $dF_p$ is the Jacobian of $\hat F$ at $\hat p$.

\bprop

    Let $F\colon M\to N$ be a map of smooth manifolds, with or without boundary, and $p\in M$.
    Then the matrix representation of $dF_p$ with respect to the coordinate bases of $M,N$ is the Jacobian of its coordinate representation $\hat F$
    at the coordinate $\hat p$ of $p$.

\eprop

Now how are two coordinate bases related?
Suppose $(U,\phi)$ and $(V,\psi)$ be two charts with $p\in U\cap V$.
Let $x^i$ be the coordinate functions of $\phi$ and $\tilde x^i$ those of $\psi$, so we can represent $v\in T_pM$ with respect to either
$\ievpdv{x^i}_p$ or $\ievpdv{\tilde x^i}_p$.

Let us denote the transition map $\psi\circ\phi^{-1}\colon\phi(U\cap V)\to\psi(U\cap V)$ by
$$ \psi\circ\phi^{-1}(x) = (\tilde x^1(x),\dots,\tilde x^n(x)) . $$
Then from our computation above,
$$ d(\psi\circ\phi^{-1})_{\phi(p)}\parens{\pdv{x^i}_{\phi(p)}} = \pdvof{\tilde x^j}{x^i}(\phi(p))\evpdv{\tilde x^j}_{\psi(p)} . $$
Now, by definition
$$ \evpdv{x^i}_p = d(\phi^{-1})_{\phi(p)}\parens{\evpdv{x^i}_{\phi(p)}} = d(\psi^{-1})_{\psi(p)}\circ d(\psi\circ\phi^{-1})_{\phi(p)}\parens{\evpdv{x^i}_{\phi(p)}}
= d(\psi^{-1})_{\psi(p)}\parens{\pdvof{\tilde x^j}{x^i}(\phi(p))\evpdv{\tilde x^j}_{\psi(p)}} , $$
and then again by the same definition,
$$ = \pdvof{\tilde x^j}{x^i}(\phi(p))\evpdv{\tilde x^j}_p . $$

So if $v=v^i\ievpdv{x^i}_p$ then
$$ v = v^i\pdv{x^i}_p = v^i\pdvof{\tilde x^j}{x^i}(\phi(p))\cdot\evpdv{\tilde x^j}_p , $$
meaning $v$'s components with respect to $\psi$ are
$$ \tilde v^j = \pdvof{\tilde x^j}{x^i}(\phi(p))\cdot v^i . $$

\bexam

    Let $\phi$ be the standard chart of $\bR^2$, and $\psi$ be the {\emphcolor polar coordinates}: it is the map on $\bR^2-0$ mapping $(r\cos\theta,r\sin\theta)$
    to $(r,\theta)$.
    So here $\tilde x_1,\tilde x_2$ in the above discussion are $r,\theta$ respectively.

    We see that
    $$ \evpdv r_p = \pdvof xr(\hat p)\evpdv x_p + \pdvof yr(\hat p)\evpdv y_p,\qquad \evpdv\theta_p = \pdvof x\theta(\hat p)\evpdv x_p + \pdvof y\theta(\hat p)\evpdv y_p . $$
    Now, $x=r\cos\theta$ so that $\pdvof xr=\cos\theta$, and so on.
    So we have
    $$ \evpdv r_p = \cos\theta\evpdv x_p + \sin\theta\evpdv y_p,\qquad \evpdv\theta_p = -r\sin\theta\evpdv x_p + r\cos\theta\evpdv y_p . $$

    So for example, if $p$'s representation in polar coordinates is $(r,\theta)=(2,\pi/2)$ we see that
    $$ \evpdv r_p = \evpdv y_p,\qquad \evpdv\theta_p = -2\evpdv x_p . $$

\eexam

\subsection{The tangent bundle}

\bdefn

    Let $M$ be a smooth manifold with or without boundary.
    Then we define $M$'s {\emphcolor tangent bundle} to be the disjoint union
    $$ TM = \coprod_{p\in M}T_pM . $$
    We denote an element in $TM$ by $(p,v)$ or $v_p$ for $p\in M$ and $v\in T_pM$.
    There is an obvious {\emphcolor projection map} $\pi\colon TM\to M$ mapping $(p,v)\mapsto p$.

\edefn

When $M=\bR^n$ we see that quite naturally
$$ T\bR^n = \coprod_{a\in\bR^n}T_a\bR^n \cong \coprod_{a\in\bR^n}\bR^n_a = \coprod_{a\in\bR^n}\set a\times\bR^n = \bR^n\times\bR^n . $$
Obviously by similar considerations in general $TM$ is in bijection with $M\times\bR^n$, but in general this does not lift to an isomorphism with richer structure;
as we will see now $TM$ has a natural topology making it into a smooth manifold.

\bprop

    For a smooth $n$-manifold $M$, $TM$ has a natural topology making it into a smooth manifold of dimension $2n$ and for which
    the projection $\pi\colon TM\to M$ is smooth.
    The boundary of $TM$ is
    $$ \partial TM = \coprod_{p\in\partial M}T_pM . $$

\eprop

\bproof

    For a chart $(U,\phi)$ of $M$, notice that $\pi^{-1}(U)\subseteq TM$ is the collection of all tangent vectors tangent to points in $U$.
    Let $x^1,\dots,x^n$ be the coordinate functions of $\phi$, and define $\tilde\phi\colon\pi^{-1}(U)\to\bR^{2n}$ by
    $$ \tilde\phi\parens{v^i\evpdv{x^i}_p} = (x^1(p),\dots,x^n(p),v^1,\dots,v^n) , $$
    the vector whose first $n$ components are $\phi(p)$, and then $v^1,\dots,v^n$.
    The image of $\tilde\phi$ is $\phi(U)\times\bR^n$, which is open in $\bR^{2n}$, and $\tilde\phi$ is a bijection onto this: we can explicitly
    write its inverse as
    $$ \tilde\phi^{-1}(x^1,\dots,x^n,v^1,\dots,v^n) = v^i\evpdv{x^i}_{\phi^{-1}(\bar x)} . $$

    Now suppose we have two smooth charts $(U,\phi),(V,\psi)$ with corresponding charts $(\pi^{-1}(U),\tilde\phi),(\pi^{-1}(V),\tilde\psi)$ in $TM$.
    Then we have
    $$ \tilde\phi(\pi^{-1}(U)\cap\pi^{-1}(V)) = \phi(U\cap V)\times\bR^n , $$
    which is open in $\bR^{2n}$, and similar for $\tilde\psi$.
    And the transition map $\tilde\psi\circ\tilde\phi^{-1}\colon\phi(U\cap V)\times\bR^n\to\psi(U\cap V)\times\bR^n$ can be computed explicitly:
    $$ \tilde\psi\circ\tilde\phi^{-1}(x^1,\dots,x^n,v^1,\dots,v^n) = \tilde\psi\parens{v^i\evpdv{x^i}_{\phi^{-1}(x^1,\dots,x^n)}} =
    (\tilde x^1,\dots,\tilde x^n,\tilde v^1,\dots,\tilde v^n) . $$
    Where $\tilde x=\psi\circ\phi^{-1}(x)$ is smooth, and $\tilde v^i$ are given as above,
    $$ \tilde v^i = \pdvof{\tilde x^i}{x^j}(x)v^j = \pdvof{\pi_i\circ\psi\circ\phi^{-1}}{x^j}v^j , $$
    which is smooth as well.

    A countable cover $\set{U_i}$ of $M$ by coordinate domains covers $TM$ by $\set{\pi^{-1}U_i}$.
    For the Hausdorff condition, for $v_p,u_q$, either $p$ and $q$ are in the same coordinate domain $U$, in which case $v_p,u_q\in\pi^{-1}U$,
    or they have disjoint coordinate domains $p\in U,q\in V$ so that $v_p\in\pi^{-1}U,u_q\in\pi^{-1}V$ are disjoint as well.
    Thus this construction endows $TM$ with a smooth structure, of dimension $2n$.

    $\pi$ is then smooth, as
    $$ \psi\circ\pi\circ\tilde\phi^{-1}(\bar x,\bar v) = \psi\circ\phi^{-1}(\bar x) $$
    is smooth.

    If $p\in\partial M$ then its boundary chart $(U,\phi)$ maps vectors in $T_pM$ to vectors in $\phi(p)\times\bR^n$, on the boundary.
    All other points are mapped to the interior.
    \qed

\eproof

The tangent bundle is functorial: for $F\colon M\to N$ smooth, we define its {\emphcolor global differential} $dF\colon TM\to TN$ to be the map whose
restriction to $T_pM\subseteq TM$ is $dF_p$; it maps $v_p\in TM$ to $dF_p(v)$.

\bprop

    The global differential is smooth, and functorial.

\eprop

\bproof

    Let $(U,\phi)$ be a chart for $M$ and $(V,\psi)$ a chart for $N$.
    We need to show that $\tilde\psi\circ dF\circ\tilde\phi^{-1}$ is smooth: let $\phi(p)=\bar x$, then
    $$ \eqalign{
        \tilde\psi\circ dF\circ\tilde\phi^{-1}(x^1,\dots,x^n,v^1,\dots,v^n) &= \tilde\psi\circ dF\parens{v^i\evpdv{x^i}_p}\cr
            &= \tilde\psi\parens{v^i\pdvof{\hat F^j}{x^i}(\psi(p))\evpdv{\tilde x^j}_{F(p)}}\cr
            &= \parens{\hat F^1(x),\dots,\hat F^n(x),v^i\pdvof{\hat F^1}{x^i}(\tilde x),\dots,v^i\pdvof{\hat F^n}{x^i}(\tilde x)},
    } $$
    where $\hat F=\psi\circ F\circ\phi^{-1}$ is the coordinate representation, and $\tilde x=\psi\circ\phi^{-1}(x)$ is the transition function.
    This is smooth, as desired.

    Functorality is clear from the functorality of local differentials.
    \qed

\eproof

Now notice that if $M$ has a global chart $\phi$, then the construction of $TM$ gives us a global chart $\tilde\phi\colon TM\to\phi(M)\times\bR^n$.
Thus $TM$ is diffeomorphic to $\phi(M)\times\bR^n$, which is diffeomorphic to $M\times\bR^n$.
So we do have that $TM\cong M\times\bR^n$ when $M$ can be covered with a global chart (i.e.\ is diffeomorphic to $\bR^n$ or $\bH^n$).
In general we will see that this is not true.

\subsection{Velocity vectors of curves}

\bdefn

    A {\emphcolor curve} on a smooth manifold $M$ is a smooth map $\gamma\colon J\to M$, from an interval $J\subseteq\bR$.
    For a point $t_0\in J$, we define $\gamma$'s {\emphcolor velocity} at $t_0$ to be
    $$ \gamma'(t_0) = d\gamma\parens{\eval{\frac d{dt}}_{t_0}} \in T_{\gamma(t_0)}(M) , $$
    where $d/dt\in T_{t_0}(\bR)$ is the standard basis vector.

\edefn

Spelled out, $\gamma'(t_0)$ is the derivation given by
$$ \gamma'(t_0)(f) = \frac{d(f\circ\gamma)}{dt}(t_0) = (f\circ\gamma)'(t_0) . $$
If $t_0\in J$ is an endpoint, this derivative is taken as one-sided.
Now let $(U,\phi)$ be a chart containing $\gamma(t_0)$, so that $\phi\gamma(t_0)=t_0$.
Then we see
$$ \gamma'(t_0) = d\gamma\parens{\eval{\frac d{dt}}_{t_0}} = \pdvof{\hat\gamma^i}{t}(t_0)\evpdv{x^i}_{\gamma(t_0)} , $$
where $\hat\gamma=\phi\circ\gamma$ is its coordinate representation, with projections $\hat\gamma^i=\pi_i\circ\hat\gamma$.
Thus $\gamma'(t_0)$ is given by essentially the same formula as in Euclidean space: by taking the partial derivatives of each component.

\bprop

    Let $M$ be a smooth manifold, with or without boundary, and $p\in M$.
    Then every $v\in T_pM$ is the velocity of some smooth curve in $M$.

\eprop

\bproof

    For an interior point $p$, let $(U,\phi)$ be an interior chart centered at $p$.
    For $v\in T_pM$, let $v=v^i\ievpdv{x^i}_p$ be its coordinate representation with respect to $\phi$, let $\bar v=(v^1,\dots,v^n)$.
    Then define
    $$ \gamma\colon(-\epsilon,\epsilon)\to M,\qquad \gamma(t)=\phi^{-1}(t\bar v) $$
    for $\epsilon>0$ small enough so that $t\bar v\in\phi(U)$ for all $t$ in the domain.
    This is a smooth curve with $\gamma(0)=p$, and its velocity at $0$ is
    $$ \gamma'(0) = \pdvof{\phi\gamma^i}{t}(0)\evpdv{x^i}_p = \pdvof{tv^i}t(0)\evpdv{x^i}_p = v^i\evpdv{x^i}_p = v , $$
    as desired.

    If $p\in\partial M$ then take a smooth chart as before and consider its coordinate representation again.
    The above formula is only defined when $tv^n\geq0$, so we take the appropriate subinterval of $(-\epsilon,\epsilon)$ so that $\gamma$ is
    defined, and we get the same result.
    \qed

\eproof

\bprop

    Let $F\colon M\to N$ be smooth, $\gamma\colon J\to M$ a smooth curve.
    Then for $t_0\in J$, the velocity of the composite curve $F\circ\gamma$ is given by
    $$ (F\circ\gamma)'(t_0) = dF(\gamma'(t_0)) . $$

\eprop

\bproof

    By definition and functorality,
    $$ (F\circ\gamma)'(t_0) = d(F\circ\gamma)\evpdv{t}_{t_0} = dF\circ d\gamma\evpdv t_{t_0} = dF(\gamma'(t_0)) . \qed $$

\eproof

While this may seem like a good way to compute the velocity vectors of composite curves, it turns out that in general it is more helpful
to view this proposition as a method of computing differentials:

\bcoro

    Let $F\colon M\to N$ be smooth, $p\in M$, and $v\in T_pM$.
    Then for any curve $\gamma\colon J\to M$ where $0\in J$ and whose velocity at $0$ is $v$, we have that
    $$ dF_p(v) = (F\circ\gamma)'(0) . $$

\ecoro

